% Exact LaTeX recreation of the supplied Phase-1.pdf
% Compile with pdfLaTeX. Place member pictures in the same folder if required.

\documentclass[11pt,a4paper]{article}

\usepackage[a4paper,left=2.0cm,right=2.0cm,top=1.55cm,bottom=1.55cm]{geometry}
\usepackage[T1]{fontenc}
\usepackage{lmodern}
\usepackage{graphicx}
\usepackage{xcolor}
\usepackage{array}
\usepackage{tabularx}
\usepackage{ragged2e}
\usepackage{enumitem}
\usepackage{fancyhdr}
\usepackage{tikz}
\usepackage{setspace}

% -------------------------------------------------
% COLOURS
% -------------------------------------------------
\definecolor{TitleBlue}{RGB}{61,103,154}

% -------------------------------------------------
% GENERAL SETTINGS
% -------------------------------------------------
\setlength{\parindent}{0pt}
\setlength{\parskip}{0pt}
\renewcommand{\arraystretch}{1.0}

% The supplied document uses a clean sans-serif heading style
% and italic serif-style explanatory text.
\newcommand{\blueheading}[1]{%
    {\sffamily\color{TitleBlue}\fontsize{15}{18}\selectfont #1}%
}

\newcommand{\bluebigheading}[1]{%
    {\sffamily\bfseries\color{TitleBlue}\fontsize{16}{19}\selectfont #1}%
}

\newcommand{\instruction}[1]{%
    {\itshape\fontsize{10.5}{12.5}\selectfont #1}%
}

\newcommand{\answerbox}[2][3.0cm]{%
    \vspace{2pt}
    \noindent
    \fbox{%
        \begin{minipage}[t][#1][t]{\dimexpr\textwidth-2\fboxsep-2\fboxrule\relax}
            \vspace{1mm}
            \itshape #2
        \end{minipage}
    }
}

% -------------------------------------------------
% HEADER / FOOTER
% -------------------------------------------------
\pagestyle{fancy}
\fancyhf{}
\renewcommand{\headrulewidth}{0pt}
\renewcommand{\footrulewidth}{0pt}

\fancyhead[L]{%
    \sffamily\bfseries\itshape\fontsize{10.5}{12}\selectfont
    BTECH CSE 
}
\fancyhead[R]{%
    \sffamily\bfseries\fontsize{10.5}{12}\selectfont
    PROJECT PROPOSAL
}
\fancyfoot[R]{%
    \itshape\fontsize{10}{12}\selectfont
    Page \thepage
}
\setlength{\headheight}{14pt}
\setlength{\headsep}{23pt}
\setlength{\footskip}{24pt}

% -------------------------------------------------
% PAGE 1
% -------------------------------------------------
\begin{document}

\vspace*{4mm}

\bluebigheading{PROJECT AND TEAM INFORMATION}

\vspace{13mm}

\blueheading{Project Title}

\vspace{1mm}

\answerbox[1.0cm]{Adaptive Network Routing & Failure Recovery Simulator}

\vspace{13mm}

\blueheading{Student / Team Information}

\vspace{2mm}

% Team information table
% Compact 4-member layout so all members remain on Page 1.
\noindent
\begin{tabularx}{\textwidth}{|>{\RaggedRight\arraybackslash}p{0.69\textwidth}|>{\RaggedRight\arraybackslash}X|}
\hline
\begin{minipage}[t]{\linewidth}
\vspace{1.5mm}
\textit{Team Name:}\\[-1mm]
\textit{Team \#015}
\vspace{1.5mm}
\end{minipage}
&
\begin{minipage}[t]{\linewidth}
\vspace{1.5mm}
\textit{Glory Hunters}
\vspace{1.5mm}
\end{minipage}
\\
\hline
\begin{minipage}[t][3.25cm][t]{\linewidth}
\vspace{1.5mm}
\textbf{Team member 1 (Team Lead)}\\[-1mm]
\vspace{2.0cm}
\end{minipage}
&
\begin{minipage}[t][3.25cm][t]{\linewidth}
\vspace{1.5mm}
\textit{Singh, Utkarsh}\\
\textit{2510012158}
\textit{2510012158@geu.ac.com}

\vspace{2mm}
\includegraphics[width=3.0cm,height=1.5cm,keepaspectratio]{image1.jpeg}
\end{minipage}
\\
\hline
\begin{minipage}[t][3.25cm][t]{\linewidth}
\vspace{1.5mm}
\textbf{Team member 2}\\[-1mm]
\vspace{2.0cm}
\end{minipage}
&
\begin{minipage}[t][3.25cm][t]{\linewidth}
\vspace{1.5mm}
\textit{Kukreti, Naitik}\\
\textit{2510011904}
\textit{2510011904@geu.ac.in}

\vspace{2mm}
\includegraphics[width=3.0cm,height=1.5cm,keepaspectratio]{image2.jpeg}
\end{minipage}
\\
\hline
\begin{minipage}[t][3.25cm][t]{\linewidth}
\vspace{1.5mm}
\textbf{Team member 3}\\[-1mm]
\vspace{2.0cm}
\end{minipage}
&
\begin{minipage}[t][3.25cm][t]{\linewidth}
\vspace{1.5mm}
\textit{Pandey, Shubh}\\
\textit{2510012661}
\textit{2510012661@geu.ac.in}
\includegraphics[width=3.0cm,height=1.5cm,keepaspectratio]{image4.jpeg}

\vspace{2mm}
\end{minipage}
\\
\hline
\begin{minipage}[t][3.25cm][t]{\linewidth}
\vspace{1.5mm}
\textbf{Team member 4}\\[-1mm]
\vspace{2.0cm}
\end{minipage}
&
\begin{minipage}[t][3.25cm][t]{\linewidth}
\vspace{1.5mm}
\textit{Bhatt, Priyansh}\\
\textit{2510012333}
\textit{2510012333@geu.ac.in}

\vspace{2mm}
 \includegraphics[width=3.0cm,height=1.5cm,keepaspectratio]{image3.jpeg}
\end{minipage}
\\
\hline
\end{tabularx}

\newpage

% -------------------------------------------------
% PAGE 2
% -------------------------------------------------

\bluebigheading{PROPOSAL DESCRIPTION}

\vspace{12mm}

\blueheading{Motivation}

\vspace{1mm}

\answerbox[6.0cm]{As we know in modern computing systems depends on networks for data transmission, that network communication is very critical. But networks are not always stable and steady. Communications links and routers can become unavailable due to failures midway or due to changing networks conditions. And this failure can affect the packet transmission , when a component on packet’s route fails packets may be delayed , lost or unable to reach their final destinations. This change in network conditions can affect and make changes during operations. A route that initially works well  can become unstable after a failure or increase in load. Therefore reliability becomes an important routing concern and shouldn’t only be considered in terms of finding a pathway. Maintaining communication when there is a change is also important. Our effective failure handling can improve continuity of communication and network performance. This motivated us to study how a network can respond to failures and adapt its routing decisions in a simulation environment. }

\vspace{12mm}

\blueheading{State of the Art / Current solution}

\vspace{1mm}

\answerbox[5.0cm]{For conventional routing , networks use routing algorithms to determine paths between source and its destination. These routing decisions can consider methods such as hop counting , link cost , latency or bandwidth. Algorithms like Dijkstra’s are commonly used as conceptual base for finding minimum-cost paths. These works well when the network’s topology and link conditions are stable. But when a router fails , routing systems can detect the failure and reroute traffic through another available path. These some of the approaches maintain alternate paths for faster recovery. More advanced systems can consider changing network conditions such as congestion , link status and topology when making rerouting decisions. Different recovery mechanisms involves clashing of things like recovery speed ,  route efficiency , congestion ,  and computational requirements like processing cost and time. This provides the basis for exploring an adaptive recovery strategy in our Stimulator.}

\vspace{12mm}

\blueheading{Project Goals and Milestones}

\vspace{1mm}

\answerbox[7.0cm]{Goals :-

1-We are building a network simulation, as we will create a model of routers , links , packets and network topology. 
2-We will implement routing by path finding algorithms and establish a base for comparison. 
3-We will stimulate network failures by introducing router failures and changing network conditions. 
4-We will implement failure recovery by detecting failures and enabling packets to use alternate routes.
5-We will develops adaptive recovery components so that system can select an appropriate recovery strategy depending on network situation and conditions.
6- finally evaluating the system through comparing different routing strategies using parameters like packet delivery , loss , recovery time , path cost , computational cost.

Milestones :-

1-Designing a network model.
2-Implement graphs and data structures.
3-Building router/link/packet/network classes.
4-Implementing routing algorithms.
5-Adding packet transmission , congestion , failure simulation , and detection.
6- Implementing adaptive recovery strategy.
7-Running experiments and comparing results.
}


% -------------------------------------------------
% PAGE 3
% -------------------------------------------------

\vspace{12mm}

\blueheading{Project Approach}

\vspace{1mm}

\answerbox[5.0cm]{We are developing project on VS code environment using C/C++ while also building UI.
Representing routers and links using graphs, using linked list to represent connections and storing properties such as link cost , status , latency , etc. Implementing Data Structure topics like Graphs for representing network , Linked Lists to maintain dynamically changing collections like paths, Queues for efficient selection of the next minimum-cost node in Dijkstra. Hash tables for fast processing of routers,  packets or links. C++ OOPS designing for modelling system using classes like router , links , packets by using features like encapsulation, inheritance, polymorphism. Stimulating link/router failures while detecting whether the current path is affected and providing alternative routes to implement adaptive recovery on them.
}


\blueheading{System Architecture (High Level Diagram)}

\vspace{1mm}

\answerbox[6.0cm]{\includegraphics[width=9cm,height=9.5cm,keepaspectratio]{image5.jpeg}}

\vspace{12mm}

\blueheading{Project Outcome / Deliverables}

\vspace{1mm}

\answerbox[5cm]{As outcome we would have built a working network simulator.  A functional program capable of modeling routers, links, packets and network topology. Having a routing functionality by possessing ability to find and use routes between source and destination nodes and also the ability to stimulate link failures ,observing their effect on packet transmission. Our system will detect affected routes and attempt to restore packet transmission through alternative routes. We will display a working implementation of adaptive recovery strategy selection for different failure situations. As experimental results comparing routing approaches using packet loss, recovery time , path cost , hop count , processing cost .}

\vspace{12mm}
\newpage

\blueheading{Assumptions}

\vspace{6mm}

\answerbox[5cm]{1-Network topology will be predefined , it will work with a defined set of routers and links rather the entire real world internet. 

2-Each link has measurable properties like cost, latency , bandwidth ; represented numerically.

3-Failures will also be simulated rather being caused by actual hardware/network failures.

4-Packets will have defined and valid source and destination in network.

5-Network changes will occur during the simulation, links/routers status ,load may change while packets are being transmitted.

6-It’s a simple model and will not attempt to reproduce every factor present in a real network ; only limited and proposed factors in project.
}

\vspace{12mm}

\blueheading{References}

\vspace{1mm}

\answerbox[3cm]{GeeksforGeeks — Graph Data Structure :- https://www.geeksforgeeks.org/dsa/graph-data-structure-and-algorithms/

GeeksforGeeks — Dijkstra's Algorithm :- https://www.geeksforgeeks.org/dsa/dijkstras-shortest-path-algorithm-greedy-algo-7/

RFC 2328 — OSPF Version 2 :- https://www.rfc-editor.org/info/rfc2328/

RFC 5880 — Bidirectional Forwarding Detection (BFD) :- https://www.rfc-editor.org/info/rfc5880/

}

\end{document}
